%
% 15-geometrisch.tex -- Geometrische Beschreibung der Rechenoperationen in C
%
% (c) 2026 Prof Dr Andreas Müller
%
\subsection{Geometrische Beschreibung der Rechenoperationen in $\mathbb{C}$}
\input{chapters/010-komplex/fig/fig-gaussebene.tex}%
Eine komplexe Zahl $z=a+bi$ kann sowohl als Vektor
\[
z = \begin{pmatrix}a\\b\end{pmatrix}
\]
in der gaussschen Zahlenebene der
\index{Zahlenebene}%
\index{Zahlenebene!gausssche}%
Abbildung~\ref{buch:komplex:zahlen:fig:gaussebene}
als auch als $2\times 2$-Matrix beschrieben werden.
Schreiben wir das Produkt $(a+bi)(c+di)$ in Matrix-Form als
\begin{align*}
(aI+bJ)(cI+dJ)
&=
\begin{pmatrix*}[r]
a&-b\\
b& a
\end{pmatrix*}
\begin{pmatrix*}[r]
c&-d\\
d& c
\end{pmatrix*}
\intertext{und betrachten wir nur die erste Spalte des Resultats,
finden wir das Produkt als Spaltenvektor}
\begin{pmatrix*}
ac-bd\\
bc+ad
\end{pmatrix*}
&=
\begin{pmatrix*}[r]
a&-b\\
b& a
\end{pmatrix*}
\begin{pmatrix*}[r]
c\\
d
\end{pmatrix*}
\end{align*}
Wie kann man dieses Produkt geometrisch in der gaussschen Zahlenebene
beschreiben?

Wir schreiben $r=\!\sqrt{a^2+b^2}$ und schreiben die Matrix von $a+bi$
als
\[
aI+bJ
=
r
\begin{pmatrix*}[r]
a/r & -b/r \\
b/r &  a/r
\end{pmatrix*}.
\]
Für die Matrixelemente gilt
\[
\biggl(\frac{a}{r}\biggr)^2
+
\biggl(\frac{b}{r}\biggr)^2
=
\frac{a^2}{r^2}
+
\frac{b^2}{r^2}
=
\frac{a^2+b^2}{r^2}
=
1.
\]
Es gibt daher einen Winkel $\varphi$ derart, dass
\[
\cos\varphi = \frac{a}{r}
\qquad\text{und}\qquad
\sin\varphi = \frac{b}{r}.
\]
Damit wird die Matrix zu
\[
aI+bJ
=
r
\begin{pmatrix*}[r]
\cos\varphi & -\sin\varphi \\
\sin\varphi &  \cos\varphi
\end{pmatrix*}
=
r\,D_\varphi,
\]
wobei $D_\varphi$ die Drehung der Ebene um den Winkel $\varphi$ im
Gegenuhrzeigersinn ist.
Die Multiplikation mit der komplexen Zahl $a+bi$ dreht also einen
Vektor in der gaussschen Zahlenebene um den Winkel $\varphi$ und
streckt ihn anschliessend um $r=|a+bi|$.

\begin{definition}[Polardarstellung komplexer Zahlen]
\index{Polardarstellung}%
Ist $z=a+bi$ eine komplexe Zahl, dann gibt es einen Winkel $\varphi$ derart,
dass 
\[
z
=
a + bi
=
r(\cos\varphi + i\sin\varphi)
\qquad\text{mit}\qquad
r=|z|.
\]
Der Winkel $\varphi$ heisst das \emph{Argument} $\varphi=\operatorname{arg}z$
der komplexen Zahl $z$.
\index{Argument einer komplexen Zahl}%
\index{argz@$\operatorname{arg}z$}%
\end{definition}

